\documentclass[11pt,a4paper]{article}

\usepackage[utf8]{inputenc}
\usepackage[T1]{fontenc}
\usepackage[spanish,es-noquoting]{babel}
\usepackage{amsmath,amssymb,amsthm}
\usepackage{booktabs}
\usepackage{longtable}
\usepackage[margin=2.6cm]{geometry}
\usepackage{tikz}
\usetikzlibrary{arrows.meta,positioning}
\usepackage{hyperref}

% Los identificadores del sistema llevan guiones y son largos. Sin permitir
% el corte por el guion, en linea desbordan el margen (hasta 216 pt medidos).
\newcommand{\ident}[1]{\textit{\nohyphens{#1}}}
\newcommand{\nohyphens}[1]{#1}
\usepackage{microtype}
\hypersetup{colorlinks=true,linkcolor=black,citecolor=black,urlcolor=blue}

\theoremstyle{plain}
\newtheorem{teorema}{Teorema}[section]
\newtheorem{corolario}[teorema]{Corolario}
\newtheorem{lema}[teorema]{Lema}
\theoremstyle{definition}
\newtheorem{definicion}[teorema]{Definición}
\theoremstyle{remark}
\newtheorem{observacion}[teorema]{Observación}

\newcommand{\Skills}{\mathcal{S}}
\newcommand{\Agentes}{\mathcal{A}}
\newcommand{\base}{\ast}

\title{El funtor cociente $\pi:\Skills\to\Agentes$\\[.3em]
\large Construcción, verificación formal en Lean 4 y medición\\
sobre el grafo de habilidades de \textsc{Metamatemático}}
\author{Leonardo Jiménez Martínez\\ \small en colaboración con Claude Opus 5 (Anthropic)}
\date{22 de agosto de 2026}

\sloppy
\begin{document}
\maketitle

\begin{abstract}
El sistema \textsc{Metamatemático} maneja objetos de tres clases —habilidades
expertas, agentes especializados y teorías matemáticas— que el diseño describía
como tres categorías $\mathbf{Graph}$ isomorfas entre sí. Este reporte establece
que son dos, no tres, y que no son isomorfas ni pueden serlo: las habilidades y
las teorías son literalmente el mismo grafo, y frente a los $172$ objetos de
éste la categoría de agentes tiene $15$, lo que es una obstrucción de
cardinalidad y no un defecto de implementación. Se construye en su lugar el
funtor cociente $\pi(s)=\text{categoría}(s)$, se demuestran en Lean~4 sus tres
propiedades relevantes —funtorialidad, preservación de co-conos y no
preservación de la minimalidad, esta última con contraejemplo finito— y se
verifican las tres sobre el grafo real. La medición confirma los teoremas:
$\pi$ es funtor ($172/172$ objetos, $226/226$ morfismos, $250$ composiciones sin
fallo), preserva los $18$ co-conos límite descubiertos ($18/18$) y conserva la
minimalidad sólo en $12$ de ellos. De los seis que la pierden, cinco tienen al
objeto base entre sus componentes, exactamente la configuración que el
contraejemplo aísla.
\end{abstract}

\tableofcontents

\section{Planteamiento}

El diseño de \textsc{Metamatemático} describe tres categorías $\mathbf{Graph}$
que deberían ser isomorfas entre ellas: una cuyos objetos son habilidades
expertas en cada rama matemática, otra cuyos objetos son agentes —uno por área—
y una tercera cuyos objetos son las teorías y subteorías. Todas ellas con una
estructura de Ehresmann sobre la complexificación que arrojan sus objetos y
flechas.

La primera observación de este trabajo es que las categorías son dos.

\begin{observacion}[Habilidades y teorías son el mismo grafo]
Los $162$ objetos de niveles \textsc{l1}--\textsc{l3} \emph{son} las ramas y
subramas. El objeto \guillemotleft{}teoría de anillos\guillemotright{} y la habilidad \guillemotleft{}experto en teoría de
anillos\guillemotright{} no son dos nodos relacionados por un isomorfismo: son el mismo nodo,
con el mismo identificador, en el mismo grafo. No hay dos copias que puedan ser
isomorfas.
\end{observacion}

La categoría genuinamente distinta es la de agentes. Y con ella el isomorfismo
tampoco es alcanzable:

\begin{observacion}[Obstrucción de cardinalidad]
$|\mathrm{Obj}(\Skills)| = 172$ y $|\mathrm{Obj}(\Agentes)| = 15$. Un
isomorfismo de categorías exige biyección en objetos. No existe.
\end{observacion}

Lo que sí existe, y es lo que hace falta para que la estructura de Ehresmann
baje de un nivel al otro, es una proyección
\[
  \pi(s) \;=\; \text{categoría}(s),
\]
sobreyectiva en objetos y muy lejos de inyectiva. El resto del reporte la
construye como funtor, demuestra qué preserva y mide qué ocurre en el grafo
real.

\section{Por qué $\pi$ no era un funtor}

Medido sobre el grafo del sistema ($172$ objetos, $556$ morfismos), la
proyección tal como estaba fallaba por dos motivos independientes.

\begin{center}
\begin{tabular}{lrr}
\toprule
Morfismos & Cantidad & \% \\
\midrule
Dentro de una misma categoría & $286$ & $51{,}4$ \\
\textbf{Cruzan} de categoría & $202$ & $36{,}3$ \\
Tocan un fundacional \textsc{l0} & $68$ & $12{,}2$ \\
\midrule
Total & $556$ & $100{,}0$ \\
\bottomrule
\end{tabular}
\end{center}

Los $202$ cruces no son errores. El morfismo
$\textit{commutative-algebra} \to \textit{algebraic-geometry}$ es
matemáticamente correcto, y es justamente el que produce el colímite
\[
  \mathrm{join}(\textit{functors},\,\textit{commutative-algebra})
  \;=\; \textit{algebraic-geometry}
\]
que el sistema descubre. El problema estaba en el
codominio: la categoría de agentes era \emph{discreta} —catorce nombres en una
lista, sin morfismos entre ellos— y no tenía dónde recibirlos.

Los $68$ restantes tocan una de las diez habilidades fundamentales, que no
llevan categoría matemática. Sobre ellas $\pi$ ni siquiera estaba definida.

\begin{observacion}[Dos clases de objeto en una sola lista]
De los $202$ cruces, $150$ son de tipo \textsc{translation} y casi todos apuntan
al mismo destino: $\textit{algebra}\to\textit{lean-tactics}$ ($33$),
$\textit{analysis}\to\textit{lean-tactics}$ ($20$),
$\textit{number-theory}\to\textit{lean-tactics}$ ($14$). Las categorías
\textit{lean-tactics} y \textit{proof-strategies} no son ramas de las
matemáticas al nivel de \textit{algebra} o \textit{topology}: son \emph{cómo} se
demuestra, no \emph{sobre qué}. Por eso todo apunta hacia ellas y nada sale
hacia las ramas. El sistema ya había tropezado con este hecho desde el otro
lado, al excluir \textsc{translation} del conjunto \texttt{ORDER\_MORPHISMS}
porque generaba uniones espurias.
\end{observacion}

En consecuencia, la construcción se restringe a los morfismos de orden
—\textsc{dependency} y \textsc{specialization}—, que son $226$.

\section{La construcción}

\begin{definicion}[Codominio inducido]
$\Agentes$ se define como la \emph{imagen} de $\pi$ con los morfismos
inducidos:
\[
  \mathrm{Obj}(\Agentes) = \pi\bigl(\mathrm{Obj}(\Skills)\bigr),
  \qquad
  \mathrm{Hom}(\Agentes) = \bigl\{\, \pi(f) : f\in\mathrm{Hom}(\Skills),\;
  \pi(f)\neq\mathrm{id} \,\bigr\}.
\]
\end{definicion}

Con esta definición $\pi$ es funtor \emph{por construcción} y no por suerte:
todo morfismo de $\Skills$ tiene imagen porque su imagen es, por definición, un
morfismo de $\Agentes$.

\begin{definicion}[Objeto base]
Las habilidades fundamentales se envían a un objeto $\base$ de $\Agentes$. La
alternativa —excluirlas del dominio— dejaría $\pi$ parcial, y con ella todo
enunciado sobre preservación. Se prefiere pagar el precio de $\base$ y
declararlo; la sección~\ref{sec:base} muestra cuál es ese precio.
\end{definicion}

El codominio resultante tiene $15$ objetos y $32$ morfismos. Los $130$ morfismos
intra-categoría se colapsan a identidades. Los inducidos de mayor multiplicidad:

\begin{center}
\begin{tabular}{lr}
\toprule
Morfismo inducido & Multiplicidad \\
\midrule
$\base \to \textit{lean-tactics}$ & $18$ \\
$\textit{lean-tactics} \to \textit{proof-strategies}$ & $13$ \\
$\base \to \textit{category-theory}$ & $8$ \\
$\textit{algebra} \to \textit{number-theory}$ & $4$ \\
$\textit{analysis} \to \textit{optimization}$ & $4$ \\
$\base \to \textit{algebra}$ & $4$ \\
$\textit{analysis} \to \textit{geometry}$ & $3$ \\
$\textit{algebra} \to \textit{topology}$ & $3$ \\
\bottomrule
\end{tabular}
\end{center}

\section{Resultados formales}

Los tres teoremas siguientes están demostrados en Lean~4 en el archivo
\texttt{MetamathProver/CategoryFoundations/QuotientFunctor.lean}, sin
\texttt{sorry} ni axiomas añadidos. El predicado \texttt{IsJoin} procede de
\texttt{JoinColimit.lean}, donde se establece que en un preorden equivale a la
propiedad universal del colímite.

\begin{teorema}[Funtorialidad]
\label{thm:funtor}
En categorías delgadas, una proyección monótona es un funtor: la ley de
identidades es la reflexividad y la de composición es la transitividad. No hay
nada más que comprobar, porque $\mathrm{Hom}(a,b)$ tiene a lo sumo un elemento.
\end{teorema}

\noindent\emph{En Lean:} \texttt{proyeccion\_es\_funtor}.

\begin{teorema}[Preservación de co-conos]
\label{thm:cocono}
Si $j$ es cota superior de $T\subseteq\mathrm{Obj}(\Skills)$, entonces $\pi(j)$
es cota superior de $\pi(T)$.
\end{teorema}

\begin{proof}
Inmediato de la monotonía: si $s\le j$ para todo $s\in T$, entonces
$\pi(s)\le\pi(j)$. \emph{En Lean:} \texttt{functor\_preserves\_cocone} y su
versión sobre \texttt{Finset}.
\end{proof}

La inmediatez de esta demostración es el punto: no hace falta ninguna hipótesis
sobre $\pi$ más allá de la monotonía. Es lo que permite que la estructura de
co-conos baje del nivel de habilidades al de agentes.

\begin{corolario}
Todo colímite se proyecta al menos a un co-cono.
\emph{En Lean:} \texttt{join\_maps\_to\_cocone}.
\end{corolario}

\begin{teorema}[La minimalidad no se preserva]
\label{thm:nomin}
Existen un preorden finito $\Skills$, un preorden finito $\Agentes$, una
proyección monótona $\pi$, un conjunto $T$ y un objeto $j$ tales que $j$ es el
join de $T$ en $\Skills$ pero $\pi(j)$ no es el join de $\pi(T)$ en $\Agentes$.
\end{teorema}

\noindent\emph{En Lean:} \texttt{functor\_not\_preserves\_join}, con
contraejemplo explícito y decidible.

\subsection{El contraejemplo necesita cinco objetos}

Un primer intento con tres objetos —$\base \le x \le j$— no funciona, y el
motivo es instructivo: en esa configuración $x$ es \emph{también} cota superior
de $\{\base,x\}$, de modo que el join es $x$ y no $j$. Colapsar por sí solo no
rompe la minimalidad.

Lo que la rompe es un \textbf{segundo representante}: dos objetos distintos del
dominio con la misma imagen, uno de los cuales abre en el codominio un atajo que
en el dominio no existe.

\begin{center}
\begin{tikzpicture}[>=Stealth, node distance=1.3cm and 2.2cm,
  every node/.style={font=\small}]
  \node (a)  {$a$};
  \node (a2) [below=of a] {$a'$};
  \node (b)  [below=of a2] {$b$};
  \node (j)  [right=of a]  {$j$};
  \node (m)  [right=of b]  {$m$};
  \draw[->] (a) -- (j);
  \draw[->] (b) -- (j);
  \draw[->] (a2) -- (m);
  \draw[->] (b) -- (m);
  \node[above=.25cm of a, xshift=1cm] {\textsc{dominio}};

  \begin{scope}[xshift=6.5cm]
    \node (A)  {$A$};
    \node (B)  [below=2.6cm of A] {$B$};
    \node (J)  [right=of A] {$J$};
    \node (M)  [right=of B] {$M$};
    \draw[->] (A) -- (J);
    \draw[->] (B) -- (J);
    \draw[->] (A) -- (M);
    \draw[->] (B) -- (M);
    \node[above=.25cm of A, xshift=1cm] {\textsc{codominio}};
  \end{scope}
\end{tikzpicture}
\end{center}

En el dominio, $m$ \emph{no} es cota superior de $\{a,b\}$, porque
$a\not\le m$; luego $j$ es el join. Pero $A\le M$ sí vale en el codominio,
porque $a'\le m$ y $\pi(a')=A$. Así $M$ se cuela como cota superior de $\{A,B\}$
sin ser imagen de ninguna cota superior de $\{a,b\}$, y como $J\not\le M$, la
imagen del join deja de ser mínima.

\begin{observacion}
El teorema~\ref{thm:nomin} no señala un defecto de la construcción. Dice que
$\pi$ es un funtor que preserva co-conos y no colímites, que es el
comportamiento normal de un cociente. Preservar co-conos y preservar co-conos
\emph{límite} son propiedades distintas, y confundirlas equivale a afirmar que
un cociente conserva colímites.
\end{observacion}

\begin{teorema}[Reflejo de la no-alcanzabilidad]
Si $\pi(a)\not\le\pi(b)$ en $\Agentes$, entonces $a\not\le b$ en $\Skills$.
\end{teorema}

\noindent\emph{En Lean:} \texttt{refleja\_no\_alcanzabilidad}. Es el
contrarrecíproco de la monotonía, y es la propiedad útil en la práctica:
permite descartar en el nivel barato ($15$ objetos) pares que no hace falta
examinar en el caro ($172$).

\section{Verificación sobre el grafo real}

El script \texttt{scripts/verificar\_funtor.py} construye $\pi$ sobre el grafo
que el sistema usa de verdad y comprueba las leyes. Todas las cifras de este
reporte proceden de \texttt{data/funtor\_pi.json}, generado por esa ejecución;
ninguna se escribe a mano.

\subsection{Funtorialidad}

\begin{center}
\begin{tabular}{lrl}
\toprule
Comprobación & Valor & Veredicto \\
\midrule
Objetos del dominio & $172$ & \\
Objetos sin imagen & $0$ & \textsc{ok} \\
Morfismos considerados & $226$ & \\
Morfismos con imagen & $226$ & \textsc{ok} \\
Composiciones comprobadas & $250$ & \\
Fallos de composición & $0$ & \textsc{ok} \\
\midrule
\multicolumn{2}{l}{\textbf{$\pi$ es funtor}} & \textbf{sí} \\
\bottomrule
\end{tabular}
\end{center}

La ley (F2) se comprueba sobre la clausura transitiva porque el codominio es un
preorden: la composición de dos flechas es la flecha compuesta, y basta con que
exista un camino.

\subsection{Preservación de los colímites descubiertos}

El motor de complexificación descubre $18$ co-conos límite y registra $27$
huecos conceptuales, con orden de complejidad máximo $\mathrm{cn}=2$. Sobre esos
$18$:

\begin{center}
\begin{tabular}{lrl}
\toprule
Propiedad & Preservados & Corresponde a \\
\midrule
Co-cono & $18/18$ & Teorema~\ref{thm:cocono} \\
Colímite (con minimalidad) & $12/18$ & Teorema~\ref{thm:nomin} \\
Colapsados a un punto & $1/18$ & --- \\
\bottomrule
\end{tabular}
\end{center}

La medición confirma los dos teoremas por separado y en la dirección esperada:
el co-cono sobrevive siempre, la minimalidad no.

\subsection{Los seis casos que pierden minimalidad}
\label{sec:base}

\begin{center}
\small
\begin{tabular}{llc}
\toprule
Colímite en $\Skills$ & Imagen por $\pi$ & $\base$ \\
\midrule
$\mathrm{join}(\textit{point-set-topology},\textit{ordinals})$
  & $\{\textit{topology},\base\}\to\textit{set-theory}$ & sí \\
$\mathrm{join}(\textit{functors},\textit{commutative-algebra})$
  & $\{\base,\textit{algebra}\}\to\textit{geometry}$ & sí \\
$\mathrm{join}(\textit{homotopy-theory},\textit{cic})$
  & $\{\textit{topology},\base\}\to\textit{logic}$ & sí \\
$\mathrm{join}(\textit{fol-deduction},\textit{tactic-exact})$
  & $\{\base,\textit{lean-tactics}\}\to\textit{proof-strategies}$ & sí \\
$\mathrm{join}(\textit{limits},\textit{exact-sequences})$
  & $\{\base,\textit{algebra}\}\to\textit{category-theory}$ & sí \\
$\mathrm{join}(\textit{tactic-apply},\textit{tactic-rewrite})$
  & $\{\textit{lean-tactics}\}\to\textit{proof-strategies}$ & no \\
\bottomrule
\end{tabular}
\end{center}

\begin{observacion}[El precio del objeto base]
Cinco de los seis tienen $\base$ entre sus componentes. La razón es
estructural y coincide con el contraejemplo: $\base$ hace de $a'$. Recibe
morfismos de habilidades fundamentales distintas y por eso queda por debajo de
casi todo el codominio —$18$ flechas hacia \textit{lean-tactics}, $8$ hacia
\textit{category-theory}, $4$ hacia \textit{algebra}, entre otras—, de modo que
el conjunto de cotas superiores de $\{\base, x\}$ es grande y la imagen del join
rara vez es la menor.

Es decir: el objeto que se introdujo para que $\pi$ fuera total es exactamente
lo que cuesta la minimalidad. Es una disyuntiva real, no un descuido, y se
resuelve a favor de la totalidad porque una $\pi$ parcial invalidaría el
teorema~\ref{thm:cocono}, que es el resultado que se quería.
\end{observacion}

El sexto caso tiene otra forma: ambas componentes colapsan a
\textit{lean-tactics} y el join va a \textit{proof-strategies}. Es el mismo
fenómeno visto en la frontera entre las dos clases de objeto que conviven en la
lista de categorías.

\section{Qué queda establecido y qué no}

\paragraph{Establecido.}
$\Skills$ y $\Agentes$ no son isomorfas y no pueden serlo. Están relacionadas
por un funtor cociente $\pi$, verificado sobre el grafo real en sus dos leyes.
Ese funtor transporta co-conos del nivel de habilidades al de agentes sin
excepción —$18/18$ medido, teorema~\ref{thm:cocono} demostrado— y refleja la
no-alcanzabilidad, lo que permite descartar en el nivel de $15$ objetos.

\paragraph{No establecido.}
$\pi$ no transporta colímites: conserva la minimalidad en $12$ de $18$, y el
teorema~\ref{thm:nomin} muestra que no puede garantizarse en general. Un
enunciado del tipo \guillemotleft{}la complexificación de Ehresmann es la misma en los dos
niveles\guillemotright{} sería falso tal cual; la forma correcta es que los co-conos se
transportan y los co-conos límite no.

\paragraph{Frente abierto.}
Tres cuestiones quedan planteadas por esta construcción.

\begin{enumerate}
\item \textbf{Refinar $\base$.} Sustituir el objeto base único por uno por
pilar (\textsc{set}, \textsc{cat}, \textsc{log}, \textsc{type}) reduciría su
grado de salida y podría recuperar parte de la minimalidad perdida. Es
comprobable con la misma medición.

\item \textbf{Separar las dos clases de objeto.} \textit{lean-tactics} y
\textit{proof-strategies} no son ramas matemáticas. Extraerlas de la lista de
categorías dejaría $\Agentes$ como taxonomía genuina y eliminaría el sexto caso
anómalo.

\item \textbf{Caracterizar cuándo sí se preserva el colímite.} Los $12$ casos
que sí lo conservan deben cumplir alguna condición suficiente; identificarla
convertiría una observación empírica en un teorema.
\end{enumerate}

\section*{Reproducibilidad}
\addcontentsline{toc}{section}{Reproducibilidad}

\begin{itemize}
\item \texttt{nucleo/graph/functor.py} — la construcción.
\item \texttt{scripts/verificar\_funtor.py} — genera
      \texttt{data/funtor\_pi.json}, fuente de toda cifra de este reporte.
\item \texttt{tests/test\_funtor.py} — $11$ tests, incluida la contraparte en
      Python del contraejemplo de Lean.
\item \texttt{MetamathProver/CategoryFoundations/QuotientFunctor.lean} — los
      cuatro teoremas. Compila con \texttt{lake build}: $6544$ trabajos,
      $0$ errores, $0$ \texttt{sorry}.
\end{itemize}

\end{document}
